%Master File:lectures.tex

\lesson{Local Sequence Matching}
\vspace{-2cm}
\begin{center}
 \includegraphics[height=10cm]{smithWaterman-29}
\end{center}

\keywords{Local sequence matching, Gaps, FSA, Extreme Values}
%%%%%%%%%%%%%%%%%%%%%%% Next Slide %%%%%%%%%%%%%%%%%%%%%%%
\renewcommand{\Outline}{%
\begin{slide}
\section[1]{Outline}

\begin{minipage}{12cm}
  \begin{enumerate}\squeeze
    \outlineitem{Local Sequence Alignment}{lsa}
    \outlineitem{Affine Gaps}{afg}
    \outlineitem{BLAST and FASTA}{blfa}
    \outlineitem{Significance Scores}{sigsco}
   \end{enumerate}
\end{minipage}\hfill
\begin{minipage}{10cm}
  \vfill
  \includegraphics[width=10cm]{fsa_gap}
  \vfill
\end{minipage}
\end{slide}
\addtocounter{outlineitem}{1}
}

\setcounter{outlineitem}{1}
\Outline
\toptarget{firstoutline}
%%%%%%%%%%%%%%%%%%%%%%% Next Slide %%%%%%%%%%%%%%%%%%%%%%%

\begin{slide}
\section{Local Alignment}

\begin{PauseHighLight}
  \begin{itemize}
  \item We saw how we could globally align sequences\pause
  \item But the nature of evolution is that often sections from one part
    of the genome are accidentally copied to another\pause
  \item Thus, we are often more interested in finding best matching
    sub-sequences\pause 
  \item In biology this problem is known as the Smith-Waterman problem
    after the authors that first posed the problem and gave an
    algorithm\pause
  \end{itemize}
\end{PauseHighLight}

\end{slide}

%%%%%%%%%%%%%%%%%%%%%%% Next Slide %%%%%%%%%%%%%%%%%%%%%%%

\begin{slide}
\section{Dynamic Programming Solution}

\begin{PauseHighLight}
  \begin{itemize}
  \item Surprisingly, by a very small tweak to the global alignment
    problem, we can solve this problem too\pause
  \item The trick is to add the option of ignoring the match with cost
    0\pause
    \begin{align*}
      F(i,j) = \max \left\{
      \begin{matrix}
        0 \\
        F(i-1,j-1) + s(a_i, b_j)\\
        F(i-1,j) + d\\
        F(i,j-1) + d
      \end{matrix}\right.
    \end{align*}
    (we use $d=-8$ below)\pause
  \end{itemize}
\end{PauseHighLight}


\end{slide}

%%%%%%%%%%%%%%%%%%%%%%% Next Slide %%%%%%%%%%%%%%%%%%%%%%%

\begin{slide}
\section[-2]{Local Alignment}
\pause
\pb
\pauselevel{=1}
\begin{center}
  \multipdf[width=\linewidth]{smithWaterman}\pause
\end{center}
\end{slide}

%%%%%%%%%%%%%%%%%%%%%%% Next Slide %%%%%%%%%%%%%%%%%%%%%%%

\begin{slide}
\section{Difference to Global Alignment}

\begin{PauseHighLight}
  \begin{itemize}
  \item The forward algorithm is the same except we insist
    $F(i,j)\geq0$\pause
  \item We also store the location of the highest $F(i,j)$\pause
  \item In the backward algorithm we start at the best value and go back
    until we reach 0\pause
  \item It is really very easy to change a global alignment to a local
    alignment\pause
  \end{itemize}
\end{PauseHighLight}

\end{slide}

%%%%%%%%%%%%%%%%%%%%%%% Next Slide %%%%%%%%%%%%%%%%%%%%%%%

\begin{slide}
\section{Repeated and Overlapping Matches}

\begin{PauseHighLight}
  \begin{itemize}
  \item So far we have considered the single best match\pause
  \item In same cases (e.g. where we have a repeated domain or motif) we
    may have one region being repeated many times within a second
    sequence\pause
  \item We can modify our dynamic programming to allow this\pause
  \item Again we may believe that one sequence should be completely
    contained in a second, then we can modify our boundary conditions to
    find this\pause
  \item Details in \textit{Biological sequence analysis} by Durbin
    \textit{et al.}\pause
  \end{itemize}
\end{PauseHighLight}

\end{slide}

%%%%%%%%%%%%%%%%%%%%%%% Next Slide %%%%%%%%%%%%%%%%%%%%%%%
\Outline % Affine Gaps
%%%%%%%%%%%%%%%%%%%%%%% Next Slide %%%%%%%%%%%%%%%%%%%%%%%

\begin{slide}
\section{Gaps}

\begin{PauseHighLight}
  \begin{itemize}
  \item As mentioned in the last lecture insertions are usually caused
    by adding a subsequence from elsewhere\pause
  \item Deletions would result from jumping a number of bases or
    treating an exon as an intron\pause
  \item Although these are rare events, once it happens then the gap
    could be quite large\pause
  \item Up to now we have made a cost of a gap proportional to its
    length, but this seems unrealistic\pause
  \end{itemize}
\end{PauseHighLight}

\end{slide}

%%%%%%%%%%%%%%%%%%%%%%% Next Slide %%%%%%%%%%%%%%%%%%%%%%%

\begin{slide}
\section{Dumb Gap Algorithm}

\begin{PauseHighLight}
  \begin{itemize}
  \item If we assign a cost $\gamma(g)$ to a gap of size $g$ then we
    could use the dynamic programme
    \begin{align*}
      F(i,j) = \max \left\{
      \begin{matrix}
        F(i-1,j-1) + s(a_i, b_j), &\\
        F(k,j) -\gamma(i-k) & \hspace{1ex} k=0,\,\ldots,\, i-1\\
        F(i,k) -\gamma(j-k) & \hspace{1ex} k=0,\,\ldots,\, j-1\\
      \end{matrix}\right.\pause
    \end{align*}
  \item Alas for two sequences of length $n$ this is a $O(n^3)$
    algorithm which is too slow for most practical applications\pause
  \item Can we do better?\pause
  \end{itemize}
\end{PauseHighLight}

\end{slide}

%%%%%%%%%%%%%%%%%%%%%%% Next Slide %%%%%%%%%%%%%%%%%%%%%%%

\begin{slide}
\section{Affine Gap}

\begin{PauseHighLight}
  \begin{itemize}
  \item If we choose a gap cost of
    \begin{align*}
      \gamma(g) = d + (g-1)\,e\pause
    \end{align*}
    then there exists a $O(n^2)$ algorithm to do matching\pause
  \item An affine transform of a variable $x$ is of the form $a\,x +b$
    (i.e. a straight line)\pause---note that $\gamma(g)$ is an affine
    function\pause
  \item Using $d=-12$ and $e=-2$ gives a good model\pause
  \item It penalises any gap, with larger gaps being slightly more
    punished\pause
  \end{itemize}
\end{PauseHighLight}

\end{slide}

%%%%%%%%%%%%%%%%%%%%%%% Next Slide %%%%%%%%%%%%%%%%%%%%%%%

\begin{slide}
\section{Affine Gap Algorithm}

\begin{PauseHighLight}
  \begin{itemize}
  \item To efficiently cope with gaps we need to introduce three
    matrices (rather than just $F(i,j)$)\pause
    \begin{itemize}
    \item $M(i,j)$ is the best score up to $(i,j)$ given $a_i$ and $b_j$
      align\pause
    \item $I_a(i,j)$ is the best score given that $a_i$ aligns with a gap\pause
    \item $I_b(i,j)$ is the best score given that $b_j$ aligns with a gap\pause
    \end{itemize}
  \end{itemize}
\end{PauseHighLight}

\end{slide}

%%%%%%%%%%%%%%%%%%%%%%% Next Slide %%%%%%%%%%%%%%%%%%%%%%%

\begin{slide}
\section{Update Equation}

\begin{minipage}{0.58\linewidth}
 \small
  \begin{align*}
    M(i,j) &= \max \left\{
    \begin{matrix}
      M(i-1,j-1) + s(a_i,b_j) \\
      I_a(i-1,j-1) + s(a_i,b_j) \\
      I_b(i-1,j-1) + s(a_i,b_j)
    \end{matrix}\right.\\
    I_a(i,j) &= \max \left\{
    \begin{matrix}
      M(i-1,j) - d \\
      I_a(i-1,j) - e
    \end{matrix}\right.\\
    I_b(i,j) &= \max \left\{
    \begin{matrix}
      M(i,j-1) - d \\
      I_b(i,j-1) - e
    \end{matrix}\right.
  \end{align*}
\end{minipage}\hfil
\begin{minipage}{0.38\linewidth}
  \begin{center}
  \includegraphics[width=\linewidth]{fsa_gap}\\
  Finite State Automaton
  \end{center}
\end{minipage}

\end{slide}

%%%%%%%%%%%%%%%%%%%%%%% Next Slide %%%%%%%%%%%%%%%%%%%%%%%

\begin{slide}
\section[-2]{Matching with Gaps}
\pause
\pauselevel{=1}
\pb
\begin{center}
  \multipdf[width=0.9\linewidth]{gap_matching}\pause
\end{center}
\end{slide}

%%%%%%%%%%%%%%%%%%%%%%% Next Slide %%%%%%%%%%%%%%%%%%%%%%%

\begin{slide}
\section{Different FSA}

\begin{PauseHighLight}
  \begin{itemize}
  \item The finite state automaton (FSA) description provides an elegant
    visualisation of dynamic programming rules\pause
  \item It makes it much easier to come up with different models
    that perhaps can better capture the evolutionary process\pause
  \item But it comes at a cost that for each new state you need to
    generate an extra matrix---slowing down the matching\pause
  \item In fact, for affine gaps we can come up with an approximate two
    state model that nearly always finds the same solution as the three
    state model (see Durbin et \textit{al.})\pause
  \end{itemize}
\end{PauseHighLight}

\end{slide}


%%%%%%%%%%%%%%%%%%%%%%% Next Slide %%%%%%%%%%%%%%%%%%%%%%%
\Outline % Blast and Fasta
%%%%%%%%%%%%%%%%%%%%%%% Next Slide %%%%%%%%%%%%%%%%%%%%%%%

\begin{slide}
\section{Inexact Matching in Practice}

\begin{PauseHighLight}
  \begin{itemize}
  \item Proteins are a few 100 amino acids long so $O(mn)$ is
    doable\pause
  \item That said, there are 100's of millions of proteins in the main
    protein database so it can take a while to look for matches against
    all known proteins\pause
  \item DNA is another matter.  DNA sequences are of the order
    $10^5$ up to $10^{11}$\pause
  \item Using dynamic programming on a large genome is not
    feasible\pause
  \end{itemize}
\end{PauseHighLight}

\end{slide}

%%%%%%%%%%%%%%%%%%%%%%% Next Slide %%%%%%%%%%%%%%%%%%%%%%%

\begin{slide}
\section{Approximate Matching}

\begin{PauseHighLight}
  \begin{itemize}
  \item Various techniques have been developed to find approximate
    matches\pause
  \item These are not guaranteed to find the best inexact match\pause
  \item They use the heuristic of searching for sections with some exact
    matches\pause
  \item Once these sections are identified they are extended and
    combined to find good matches that are polished using Smith-Waterman\pause
  \item These algorithms are a rag-bag of heuristic techniques that are
    very effective if not very elegant\pause
  \end{itemize}
\end{PauseHighLight}

\end{slide}

%%%%%%%%%%%%%%%%%%%%%%% Next Slide %%%%%%%%%%%%%%%%%%%%%%%

\begin{slide}
\section[-2]{FASTA}

\begin{PauseHighLight}
  \begin{itemize}
  \item FASTA is an approximate matching algorithm for both proteins and
    nucleotide sequences\pause
  \item It firstly checks for short ($k$-mer) exact matches between a query sequence and
    a target\pause{} ($k =2$ for proteins and $k=4$--$6$ for DNA)\pause
  \item Local regions are identified with a high density of $k$-mers\pause
  \item Local regions are extended by looking for inexact matches\pause
  \item Ungapped regions are combined into gapped region\pause
  \item Finally plausible regions are rematched using dynamical programming\pause
  \end{itemize}
\end{PauseHighLight}

\end{slide}


%%%%%%%%%%%%%%%%%%%%%%% Next Slide %%%%%%%%%%%%%%%%%%%%%%%

\begin{slide}
\section[-2]{BLAST}

\begin{PauseHighLight}
  \begin{itemize}
  \item BLAST (Basic Local Alignment Search Tool) is the most commonly
    used tool for alignment (mainly superseding FASTA)\pause
  \item After removing low-complexity regions that might give high
    matches, exact matches (slightly larger than FASTA) are found\pause
  \item BLAST builds up high scoring ``words'' from the exact matches\pause
  \item These seeds are extended into regions which can then be combined
    (in BLAST2) to find gapped regions\pause
  \item Again these are polished with Smith-Waterman\pause
  \item For detail see Wikipedia or BLAST homepages\pause
  \end{itemize}
\end{PauseHighLight}

\end{slide}


%%%%%%%%%%%%%%%%%%%%%%% Next Slide %%%%%%%%%%%%%%%%%%%%%%%
\Outline % Significance Scores
%%%%%%%%%%%%%%%%%%%%%%% Next Slide %%%%%%%%%%%%%%%%%%%%%%%

\begin{slide}
\section{Are Matches Real?}

\begin{PauseHighLight}
  \begin{itemize}
  \item By chance we are likely to find some good matches even between
    random sequences\pause
  \item What we would like to know is the probability of finding a
    particular match by chance\pause
  \item This is used in algorithms like BLAST to decide which matched
    regions it should keep\pause
  \item The problem is non-trivial because we are choosing the best
    match that we found\pause
  \end{itemize}
\end{PauseHighLight}

\end{slide}

%%%%%%%%%%%%%%%%%%%%%%% Next Slide %%%%%%%%%%%%%%%%%%%%%%%

\begin{slide}
\section{Probability of a Match}

\begin{PauseHighLight}
  \begin{itemize}
  \item Computing the probability of an exact match between two random
    strings is possible if we know the frequency of occurrence of all
    the elements\pause
  \item For inexact matching this requires a bit more work\pause
  \item The BLOSUM50 scores are designed to give the (log-)probability of a
    particular match between amino acid residues\pause
  \item A lot of effort has gone into creating matching schemes that
    give the probability of a match\pause
  \item But what is the probability of the best of $k$ matches being
    greater than $Z$?\pause
  \end{itemize}
\end{PauseHighLight}

\end{slide}

%%%%%%%%%%%%%%%%%%%%%%% Next Slide %%%%%%%%%%%%%%%%%%%%%%%

\begin{slide}
\section[-2]{Extreme Values}
\pause
\pauselevel{=1}
\pb
\begin{center}
  \multipdf[width=0.9\linewidth]{evd}\pause
\end{center}
\end{slide}

%%%%%%%%%%%%%%%%%%%%%%% Next Slide %%%%%%%%%%%%%%%%%%%%%%%

\begin{slide}
\section{Extreme Value Distribution}


\begin{PauseHighLight}
  \begin{itemize}
  \item The distribution of extreme values (best of $k$) is known as the
    extreme value distribution\pause
  \item Similar to the central limit theorem (which states that the sum
    of random variables will be normally distributed), it turns out that
    asymptotically ($k\rightarrow\infty$) the extreme value of random
    deviates converges to a small number of extreme value
    distributions\pause
  \item The most common is the Gumbel EVD given by
    \begin{align*}
        f(z) \approx \frac{1}{\beta}
      \e{-\frac{x-\mu}{\beta}-e^{-\frac{x-\mu}{\beta}}}\pause
    \end{align*}
  \end{itemize}
\end{PauseHighLight}

\end{slide}

%%%%%%%%%%%%%%%%%%%%%%% Next Slide %%%%%%%%%%%%%%%%%%%%%%%

\begin{slide}
\section[-2]{Using EVD}

\begin{PauseHighLight}
  \begin{itemize}
  \item Extreme Value Distributions become more accurate when we do a
    large number of trials\pause
    \begin{center}
      \includegraphics[width=0.6\linewidth]{gumbel}
    \end{center}
  \item In matching the number of matches are exponential in the string
    length so we are in a situation where EVDs are quite accurate\pause
  \item The probability of observing a match with score at least $Z$ is
    {\small \begin{align*}
      \Prob{z\geq Z} = \int^{\infty}_Z f(z)\,\dd z = 
      1-\exp\left(-e^{-\frac{x-\mu}{\beta}}\right)\pause
    \end{align*}}
  \end{itemize}
\end{PauseHighLight}

\end{slide}

%%%%%%%%%%%%%%%%%%%%%%% Next Slide %%%%%%%%%%%%%%%%%%%%%%%

\begin{slide}
\section{Summary}

\begin{PauseHighLight}
  \begin{itemize}
  \item Dynamics programming provides great flexibility in
    matching\pause
  \item Using finite state automata we can develop more sophisticated
    models\pause
  \item However, for speed, we often sacrifice accuracy (i.e. we might
    miss some good matches)\pause
  \item We need ways to measure the significant of found matches\pause
  \end{itemize}
\end{PauseHighLight}

\end{slide}
